In the field of stuttered speech detection, there are several proposed datasets, such as UClass~\cite{Howell_uclass}, KSoF~\cite{bayerl_KSoFKasselState_2022}, and Sep-28K~\cite{lea_Sep28k}.
UClass is an unlabeled dataset while the others are labeled.
KSoF and Sep-28K share the same format.
Both of them consist of 3-second clips with utterance-level annotations.
Specifically, Sep-28K also includes a subset of utterance level annotation from Fluency Bank.
Not only the size of the dataset are relatively small, but they only provide 
% word-level labels, 
utterance-level labels,
which we believe is one of the main obstacles in this field.
% KSoF includes German stuttered speech with utterance labeled annotations.
% In~\cite{lea_Sep28k}, they provide utterance level annotations for SEP-28K and Fluency Bank, which contains total of 15.6hr.
% Not only the size of the dataset are relatively small, they only provide word level labels.

Several recent works have used Speech SSL models in stuttering detection.
In ~\cite{bayerl22b_interspeech}, they adopted Wav2vec 2.0 as their backbone and trained different models using different upstream layers of Wav2vec 2.0 for stuttering detection.
They also proposed to employ gender prediction as an auxiliary task for enhancing performance.
In ~\cite{bayerl23_interspeech}, they pointed out that stuttering types do not happen independently, so they framed the detection problem as a multi-label classification problem.
In ~\cite{Payal_iccasp_Siamese}, they improved the quality of Wav2vec 2.0 embeddings for stuttered speech detection using Siamese network and contrastive training. 
Despite these impressive results, their models predict at the utterance level, which is insufficient for clinical usage.
Furthermore, the optimal strategies for utilizing Speech SSL models for stuttering detection are still not fully understood.
Several prior works also incorporate Speech SSL models into their model design, such as for detecting dysarthria~\cite{Farhad_Dysarthria} and aphasia~\cite{Lian2023UnconstrainedDM,Lian2024TowardsHS}, but they focus on other types of speech disorders and are not directly applicable to our task.

To the best of our knowledge, ~\cite{harvill22_interspeech} is the only study that focused on fine-grained stuttering detection, more specifically frame-level stuttering detection.
Hence, we chose it as our baseline.
It is first pretrained on LibriSpeech with synthetic disfluencies augmented during training with frame-level supervision.
Then, they added a max pooling layer to finetune on utterance-level SEP-28K dataset.
% We are also aware that ~\cite{Lian2023UnconstrainedDM,Lian2024TowardsHS} also focused on fine-grained disfluency detection.
% Nonetheless, they hand crafted some heuristic to identify different disfluency types for Aphasia speech, while we are focusing on stuttered speech detection.

% The scale of dataset has been a bottleneck for stuttering detection currently, most of the previous work employ synthetic data augmentation on large scale clean speech datset for pretraining then fintune on real stuttering dataset.
% However, stuttering behavior is a intricate characteristic that is hard to micmic from simple synthetic data augmentation, resulting in a large generalization gap between pretraining and fintune.
% With the surge of Speech SSL models on all sorts of speech processing downstream tasks, recent works have switch their model backbones form LSTM, TDNN to speech SSl model such as wav2vec 2.0.




% Recent literature has proposed various SSL algorithms~\cite{hsu2021hubert,Chen2021WavLMLS,Conneau2020xlsr,schneider2019wav2vec} and many applications have been built on top of these SSL speech models~\cite{mohamed2022self}, including Automatic Speech Recognition~\cite{baevski2021wav2vec-u,liu2022wav2vec-u2}, Visually Grounded Speech~\cite{peng2022fastvgs,peng2022vghubert,shihSpeechCLIP,layne_mspchclip}, Speech Segmentation~\cite{peng2022vghubert,Strgar_2022}
% Emotion recognition~\cite{Morais_emotion_ssl}, speaker verification~\cite{Peng_SSL_sv_2023} and so on.
% Furthermore, there are benchmarks proposed to evaluate speech SSL models on various downstream tasks~\cite{yang21c_superb,tsai2022superbsg,shon-etal-2023-slue,shi23g_mlsuperb}.